% ============================================
% ЗАКЛЮЧЕНИЕ
% ============================================

В ходе настоящей работы решена задача снижения энергопотребления устройств интернета вещей с функцией голосового управления. Получены следующие основные результаты.

1. Проведён анализ существующих методов и систем распознавания ключевых слов. Рассмотрены облачные решения, специализированные аппаратные платформы (ASIC, FPGA, DSP) и микроконтроллеры общего назначения. Установлено, что существующие работы рассматривают каскадную архитектуру, квантизацию модели и декомпозицию энергопотребления изолированно, но не в совокупности.

2. Предложен каскадный метод детекции с тремя этапами возрастающей вычислительной сложности: аналоговая детекция звукового события, захват аудио с извлечением MFCC-признаков и нейросетевой инференс квантизованной модели. Метод допускает реализацию на различных микроконтроллерных платформах, поддерживающих режимы глубокого сна с внешним пробуждением. Экспериментальная апробация проведена на базе микроконтроллера ESP32-S3.

3. Исследовано влияние методов квантизации на модель DS-CNN-M. Посттренировочная квантизация (PTQ) и квантизация с дообучением (QAT) обеспечили сжатие модели более чем в 3,5 раза при потере точности классификации менее 0,5 процентного пункта и сокращении времени инференса более чем вдвое.

4. Разработан лабораторный стенд с разделением измеряемой системы и системы измерений в независимые домены электропитания, обеспечивающий корректное измерение потребления в диапазоне от единиц микроампер до сотен миллиампер.

5. Выполнена декомпозиция энергопотребления каскадного цикла по этапам обработки. Установлено, что вычислительные этапы (извлечение признаков и инференс) составляют 30\% энергопотребления рабочего цикла, тогда как доминирующими являются захват аудио (33\%) и холодный старт микроконтроллера (29\%). Каскадный метод обеспечивает увеличение автономности в 4--9 раз по сравнению с режимом непрерывного мониторинга.

Практическая значимость работы состоит в том, что предложенный каскадный метод и методика измерений могут быть использованы при проектировании автономных IoT-устройств с голосовым управлением. Метод не привязан к конкретной аппаратной платформе и применим к микроконтроллерам семейств STM32, nRF52 и другим, поддерживающим режимы глубокого сна с внешним пробуждением.
